% GENERATED FILE. Edit canonical lessons, then run npm run generate:books.
% canonical-source-hash: fnv1a64:8a77690774f66c99
% canonical-lessons: ZH-C53-shuo, ZH-W53-shuo, ZH-C53-hua, ZH-W53-hua, ZH-C53-ting, ZH-C53-dong, ZH-C53-man

\chapter{Speaking and Hearing}
\label{ch:zh-speaking-and-hearing}
\hlchaptermodality{53}

\begin{chapteropening}
\textbf{By the end of this chapter:} \emph{I can say that I understand or not, and ask somebody to speak slowly.}

Complete the last lesson of chapter 53: I can say that I understand or not, and ask somebody to speak slowly.
\end{chapteropening}

\section[\texorpdfstring{\zh{说} (shuō)}{说 (shuō)}]{\zh{说} (shuō) — to say; to speak}

\label{lesson:ZH-C53-shuo}



\begin{quote}
Before the new one: say the Mandarin for eye, then the Mandarin for ear.
\end{quote}

\subsection*{You'll want to know: \zh{说}}
\textbf{\zh{说}} — \emph{shuō} — \textquotedblleft{}to say; to speak\textquotedblright{}.

\begin{quote}
\textbf{\zh{他说你的手很好看}} \emph{tā shuō nǐ de shǒu hěn hǎokàn} — He says you have lovely hands.
\end{quote}

The new character, in pieces: \zh{说} — \zh{讠}, the speech piece, on the left; on the right a dot and a short fall, then \zh{口}, then two legs. The next lesson puts it in your hand.

\begin{cousinweb}
\textbf{\zh{说}} (\emph{shuō}) is inherited from Middle Chinese \emph{sywet} (Guangyun, in Baxter's transcription), glossed 'to tell'.
\end{cousinweb}

\begin{grammarlens}[title={What you've built}]
One more word for asking politely.
\end{grammarlens}

\subsection*{Your turn}
\begin{itemize}
\raggedright
  \item \emph{Say it:} \emph{shuō}
  \item \emph{Say it:} \emph{shuō}, once more
  \item \emph{Say it:} \emph{tā shuō nǐ de shǒu hěn hǎokàn} — He says you have lovely hands.
  \item \emph{Recall:} say the Mandarin for eye, then the Mandarin for ear, then say \emph{shuō} again
\end{itemize}

\begin{tcolorbox}[breakable,colback=teal!4,colframe=teal!35!black,title={Before you move on}]
What does \zh{说} mean? (\textquotedblleft{}To say; to speak\textquotedblright{}.) Say it once more.
\end{tcolorbox}
\section[\texorpdfstring{\zh{说} (shuō)}{说 (shuō)}]{\zh{说} — the speech piece, then a mouth with legs}

\label{lesson:ZH-W53-shuo}



\begin{quote}
The word first, then the character inside it.

\begin{itemize}
\raggedright
  \item \emph{Recall:} say \emph{shuō} — \textbf{R1}, one lesson back, before it decays
\end{itemize}

\begin{quote}
\zh{讠} \zh{请}
\end{quote}

\emph{Read it:} each one aloud, with nothing said first — the speech piece, \emph{qǐng}
\end{quote}

\subsection*{Script — the speech piece, then a mouth with legs}
\hlblockfigure{\detokenize{figures/ZH-W53-shuo-filmstrip.pdf}}{How \zh{说} is written, stroke by stroke}

\begin{quote}
\zh{说}
\end{quote}

\textbf{Two pieces, left before right.}

\textbf{The speech piece, \zh{讠}, two strokes}, as in \zh{请}: a dot, then a short horizontal that turns down and finishes rising to the upper right, all one stroke.

\textbf{The right piece, seven strokes, top to bottom.} A dot and a short left-falling stroke, leaning in toward each other. \zh{口} under them. Then two legs: a left-falling stroke, and a vertical that bends right along the bottom and hooks up.

2 + 7 = nine, and five of them are \zh{讠} and \zh{口}.

\subsection*{Your turn}
\begin{itemize}
\raggedright
  \item \emph{Write it:} \zh{说}, all nine, with the model in front of you
  \item \emph{Cover:} write it again from memory, and count the pieces, not the strokes
\end{itemize}

\begin{tcolorbox}[breakable,colback=teal!4,colframe=teal!35!black,title={Before you move on}]
How many pieces in \zh{说}? (\textbf{Two} — \zh{讠}, then the right side.) Which character you know is in the middle of the right side? (\textbf{\zh{口}}.) How does the last stroke end? (\textbf{It hooks up}.)

Source: \href{https://raw.githubusercontent.com/chanind/hanzi-writer-data/68d10a4b21150cae5e1ebbd223eed289cf32d90c/data/\%E8\%AF\%B4.json}{Hanzi Writer Data \zh{说}, pinned stroke paths and medians}.
\end{tcolorbox}
\section[\texorpdfstring{\zh{话} (huà)}{话 (huà)}]{\zh{话} (huà) — words; speech}

\label{lesson:ZH-C53-hua}



\begin{quote}
Before the new one: say the Mandarin for ear, then the Mandarin for to say.
\end{quote}

\subsection*{You'll want to know: \zh{话}}
\textbf{\zh{话}} — \emph{huà} — \textquotedblleft{}words; speech\textquotedblright{}.

\begin{quote}
\textbf{\zh{他说中国话}} \emph{tā shuō Zhōngguó huà} — He speaks Chinese.
\end{quote}

The new character, in pieces: \zh{话} — \zh{讠} on the left; on the right a short fall, a bar and a vertical, standing on \zh{口}. The next lesson puts it in your hand.

\begin{cousinweb}
\textbf{\zh{话}} (\emph{huà}) is inherited from Middle Chinese \emph{hwaejH} (Guangyun, in Baxter's transcription), glossed 'talk'.
\end{cousinweb}

\begin{grammarlens}[title={What you've built}]
One more word for asking politely.
\end{grammarlens}

\subsection*{Your turn}
\begin{itemize}
\raggedright
  \item \emph{Say it:} \emph{huà}
  \item \emph{Say it:} \emph{huà}, once more
  \item \emph{Say it:} \emph{tā shuō Zhōngguó huà} — He speaks Chinese.
  \item \emph{Recall:} say the Mandarin for ear, then the Mandarin for to say, then say \emph{huà} again
\end{itemize}

\begin{tcolorbox}[breakable,colback=teal!4,colframe=teal!35!black,title={Before you move on}]
What does \zh{话} mean? (\textquotedblleft{}Words; speech\textquotedblright{}.) Say it once more.
\end{tcolorbox}
\section[\texorpdfstring{\zh{话} (huà)}{话 (huà)}]{\zh{话} — the speech piece, then a mouth}

\label{lesson:ZH-W53-hua}



\begin{quote}
The word first, then the character inside it.

\begin{itemize}
\raggedright
  \item \emph{Recall:} say \emph{huà} — \textbf{R1}, one lesson back, before it decays
\end{itemize}

\begin{quote}
\zh{说} \zh{语}
\end{quote}

\emph{Read it:} each one aloud, with nothing said first — \emph{shuō}, \emph{yǔ}
\end{quote}

\subsection*{Script — the speech piece, then a mouth}
\hlblockfigure{\detokenize{figures/ZH-W53-hua-filmstrip.pdf}}{How \zh{话} is written, stroke by stroke}

\begin{quote}
\zh{话}
\end{quote}

\textbf{Two pieces, left before right.}

\textbf{The speech piece, \zh{讠}, two strokes}, as in \zh{说}.

\textbf{The right piece, six strokes, top to bottom.} A short left-falling stroke along the top. A long horizontal. A vertical down from the top through the bar. Then \zh{口} at the foot.

2 + 3 + 3 = eight.

\subsection*{Your turn}
\begin{itemize}
\raggedright
  \item \emph{Write it:} \zh{话}, all eight, with the model in front of you
  \item \emph{Cover:} write it again from memory, and count the pieces, not the strokes
\end{itemize}

\begin{tcolorbox}[breakable,colback=teal!4,colframe=teal!35!black,title={Before you move on}]
How many pieces in \zh{话}? (\textbf{Two} — \zh{讠}, then the right side.) What does the right side stand on? (\textbf{\zh{口}}.) Which character did your hand write a moment ago, with the same left piece? (\textbf{\zh{说}}.)

Source: \href{https://raw.githubusercontent.com/chanind/hanzi-writer-data/68d10a4b21150cae5e1ebbd223eed289cf32d90c/data/\%E8\%AF\%9D.json}{Hanzi Writer Data \zh{话}, pinned stroke paths and medians}.
\end{tcolorbox}
\section[tīng]{tīng — to listen; to hear}

\label{lesson:ZH-C53-ting}



\begin{quote}
Before the new one: say the Mandarin for to say, then the Mandarin for spoken words.
\end{quote}

\subsection*{You'll want to know: tīng}
\emph{tīng} — \textquotedblleft{}to listen; to hear\textquotedblright{}.

\begin{quote}
\emph{nǐ tīng wǒ shuō} — Listen to me.
\end{quote}

Its characters come later; for now it is a word for the ear.

\begin{cousinweb}
\emph{tīng} is inherited from Middle Chinese \emph{theng} (Guangyun, in Baxter's transcription), glossed 'to listen'.
\end{cousinweb}

\begin{grammarlens}[title={What you've built}]
One more word for asking politely.
\end{grammarlens}

\subsection*{Your turn}
\begin{itemize}
\raggedright
  \item \emph{Say it:} \emph{tīng}
  \item \emph{Say it:} \emph{tīng}, once more
  \item \emph{Say it:} \emph{nǐ tīng wǒ shuō} — Listen to me.
  \item \emph{Recall:} say the Mandarin for to say, then the Mandarin for spoken words, then say \emph{tīng} again
\end{itemize}

\begin{tcolorbox}[breakable,colback=teal!4,colframe=teal!35!black,title={Before you move on}]
What does tīng mean? (\textquotedblleft{}To listen; to hear\textquotedblright{}.) Say it once more.
\end{tcolorbox}
\section[dǒng]{dǒng — to understand}

\label{lesson:ZH-C53-dong}



\begin{quote}
Before the new one: say the Mandarin for spoken words, then the Mandarin for to listen.
\end{quote}

\subsection*{You'll want to know: dǒng}
\emph{dǒng} — \textquotedblleft{}to understand\textquotedblright{}.

\begin{quote}
\emph{wǒ bù dǒng} — I don't understand.
\end{quote}

Its characters come later; for now it is a word for the ear.

\begin{cousinweb}
\emph{dǒng} has no older form in any dictionary we use, so its history is not settled.
\end{cousinweb}

\begin{grammarlens}[title={What you've built}]
One more word for asking politely.
\end{grammarlens}

\subsection*{Your turn}
\begin{itemize}
\raggedright
  \item \emph{Say it:} \emph{dǒng}
  \item \emph{Say it:} \emph{dǒng}, once more
  \item \emph{Say it:} \emph{wǒ bù dǒng} — I don't understand.
  \item \emph{Recall:} say the Mandarin for spoken words, then the Mandarin for to listen, then say \emph{dǒng} again
\end{itemize}

\begin{tcolorbox}[breakable,colback=teal!4,colframe=teal!35!black,title={Before you move on}]
What does dǒng mean? (\textquotedblleft{}To understand\textquotedblright{}.) Say it once more.
\end{tcolorbox}
\section[màn]{màn — slow}

\label{lesson:ZH-C53-man}



\begin{quote}
Before the new one: say the Mandarin for to listen, then the Mandarin for to understand.
\end{quote}

\subsection*{You'll want to know: màn}
\emph{màn} — \textquotedblleft{}slow\textquotedblright{}.

\begin{quote}
\emph{tā shuō huà hěn màn} — He speaks slowly.
\end{quote}

Its characters come later; for now it is a word for the ear.

\begin{cousinweb}
\emph{màn} is inherited from Middle Chinese \emph{maenH} (Guangyun, in Baxter's transcription), glossed 'lazy, haughty'. The sense 'slow' is not given there; from lazy to slow is a likely step, but it is not settled.
\end{cousinweb}

\begin{grammarlens}[title={What you've built}]
One more word for asking politely.
\end{grammarlens}

\subsection*{Your turn}
\begin{itemize}
\raggedright
  \item \emph{Say it:} \emph{màn}
  \item \emph{Say it:} \emph{màn}, once more
  \item \emph{Say it:} \emph{tā shuō huà hěn màn} — He speaks slowly.
  \item \emph{Recall:} say the Mandarin for to listen, then the Mandarin for to understand, then say \emph{màn} again
\end{itemize}

\begin{tcolorbox}[breakable,colback=teal!4,colframe=teal!35!black,title={Before you move on}]
What does màn mean? (\textquotedblleft{}Slow\textquotedblright{}.) Say it once more.
\end{tcolorbox}
